%\section{Supporting Results on CGMT}\label{SM cgmt res}
The following theorem replaces the compactness constraint with  closedness in the CGMT and is borrowed from \cite{li2020exploring}. For related statements see also \cite[App.~A]{deng2019model}.
\begin{theorem} [CGMT with Closedness Constrains]\label{thm closed} Let $\psi$ be a convex function obeying $\lim_{\tn{\w}\rightarrow\infty}\psi(\w)=\infty$. Given a closed set $\Sc$, define
\begin{align}
\Phi_\la(\X)&=\min_{\w\in\Sc}\la\tn{\X\w}+\psi(\w)\\
\phi_\la(\g,\h)&=\min_{\w\in\Sc}\la(\tn{\w}\tn{\g}-\h^T\w)_++\psi(\w),
\end{align}
and
\begin{align}
&\Phi_\infty(\X)=\min_{\w\in\Sc,\X\w=0}\psi(\w)\\
&\phi_\infty(\g,\h)=\min_{\w\in\Sc,\tn{\w}\tn{\g}\leq \h^T\w}\psi(\w).
\end{align}
For all $\la\in[0,\infty)\cup\{\infty\}$, we have that
\begin{itemize}
\item $\Pro(\Phi_\la(\X)<t)\leq2\Pro(\phi_\la(\X)\leq t)$.
\item If $\Sc$ is additionally convex, we additionally have that $\Pro(\Phi_\la(\X)>t)\leq2\Pro(\phi_\la(\X)\geq t)$. Combining with the first statement, this implies that for any $\mu,t>0$
\[
\Pro(|\Phi_\la(\X)-\mu|>t)\leq2\Pro(|\phi_\la(\X)-\mu|\geq t)
\]
\end{itemize}
\end{theorem}


%\begin{lemma} [AO solution to PO solution] \label{dist cont lem}Let $\X\in\R^{n\times p},\g\in\R^n,\h\in\R^p\distas\Nn(0,1)$. Suppose we have two loss functions $\Lc_{PO}(\w;\X)$ and $\Lc_{AO}(\w;\g,\h)$ as a function of $\w$\footnote{$\Lc(\w,\ab)$ can account for additional set constraints of type $\w\in\Cc$ by adding the indicator penalty $\max_{\la\geq 0}\la 1_{\w\not\in\Cc}$.}. Given a set $\Sc$, define the objectives 
%\begin{align}
%\Phi_{\Sc}(\X)=\min_{\w\in\Sc}\Lc_{PO}(\w;\X)\quad\text{and}\quad\phi_{\Sc}(\g,\h)=\min_{\w\in\Sc}\Lc_{AO}(\w;\g,\h).\label{general phi}
%\end{align}
%Suppose $\Phi$ and $\phi$ satisfies the following conditions for any closed set $\Sc$ and $t$
%\begin{itemize}
%\item $\Pro(\Phi_{\Sc}(\X) < t)\leq2\Pro(\phi_{\Sc}(\g,\h)\leq t)$.
%\item Furthermore, if $\Sc$ is convex, $\Pro(\Phi_{\Sc}(\X) > t)\leq2\Pro(\phi_{\Sc}(\g,\h)\geq t)$.%$\Pro(|\Phi_{\Sc}(\X)-\mu|> t)\leq2\Pro(|\phi_{\Sc}(\g,\h)-\mu|> t)$.
%\end{itemize}
%Define the set of global minima $\Mc=\{\w\bgl \Lc(\w;\X)=\Phi(\X)\}$. For any closed set $\Sc$, we have that
%\begin{align}
%\Pro(\Mc\in\Sc^c)\geq 1-2\min_{t}(\Pro(\phi_{\R^p}(\g,\h)\geq t)+\Pro(\phi_{\Sc}(\g,\h)\leq t)).\label{min inside}
%\end{align}
%\end{lemma}


\cmt{
\section{Proof of Constrained CGMT}
\subsection{Proof for the convex case}
\begin{lemma} \label{lem convex}Given a convex and compact $\Sc$, define the PO and AO problems% under interpolation as follows
\begin{align}
&\Phi_\infty(\X)=\min_{\w\in\Sc,\X\w=0}\psi(\w)\\
&\phi_\infty(\g,\h)=\min_{\w\in\Sc,\tn{\w}\tn{\g}\leq \h^T\w}\psi(\w).
\end{align}
Suppose $\X,\g,\h\distas\Nn(0,1)$. Then, we have that
\begin{align}
\Pro(\Phi_\infty(\X)> t)\leq2\Pro(\phi_\infty(\g,\h)\geq t).
\end{align}
\end{lemma}

\subsection{Proof for the general case}
\begin{lemma} \label{lem general constraint}Given a compact set $\Sc$, define the PO and AO problems as in Lemma \ref{lem convex}. We have that
\begin{align}
\Pro(\Phi_\infty(\X)<t)\leq2\Pro(\phi_\infty(\g,\h)< t).
\end{align}
\end{lemma}
\begin{proof} The proof is similar to that of Lemma \ref{lem convex}. For a general compact set $\Sc$, application of Gordon's theorem yields the one-sided bound
\begin{align}
\Pro(\Phi_\la(\X)<t)\leq2\Pro(\phi_\la(\g,\h)\leq t).
\end{align}
To move from finite $\la$ to infinite, we make use of Lemma \ref{lem continuous limit}. Define the indicator function $E_\la=1_{\Phi_\la(\X)\leq t}$. Using Lemma \ref{lem continuous limit}, for any choice of $\X$, $\lim_{\la\rightarrow\infty}E_\la=\lim_{\la\rightarrow\infty}1_{\Phi_\la(\X)< t}=1_{\Phi_\infty(\X)< t}$. Note again that, if the problem is infeasible, then $\lim_{\la\rightarrow\infty}E_\la=E_\infty=0$. To proceed, we are in a position to apply Dominated Convergence Theorem to find
\begin{align}
&\lim_{\la \rightarrow\infty}\E[E_\la]=\E[E_\infty]\iff\Pro(\Phi_\infty(\X)< t)=\lim_{\la\rightarrow\infty}\Pro(\Phi_\la(\X)< t).
\end{align}
Applying the identical argument on $\phi_{\g,\h}$ to find $\Pro(\phi_\infty(\g,\h)\leq t)=\lim_{\la\rightarrow\infty}\Pro(\phi_\la(\g,\h)\leq t)$, we obtain the desired relation
\begin{align}
\Pro(\Phi_\infty(\X)< t)&=\lim_{\la\rightarrow\infty}\Pro(\Phi_\la(\X)< t)\\
&\leq 2\lim_{\la\rightarrow\infty}\Pro(\phi_\la(\g,\h)\leq t)\\
&= 2\Pro(\phi_\infty(\g,\h)\leq t).
\end{align}
\end{proof}
\begin{lemma}\label{lem continuous limit} Let $\Sc$ be a compact set and $\psi(\cdot)$ be a continuous function and $f(\w)$ be a non-negative continuous function. Then
\[
\lim_{\la\rightarrow\infty}\min_{\w\in\Sc}\la f(\w)+\psi(\w)=\min_{\w\in\Sc,f(\w)=0}\psi(\w)
\]
Thus, setting $f(\w)=\tn{\X\w}$ and $f(\w)=\tn{\w}\tn{\g}-\h^T\w$, we have that
\begin{align*}
&\lim_{\la\rightarrow\infty}\Phi_{\la}(\X)=\Phi_{\infty}(\X)\\
&\lim_{\la\rightarrow\infty}\phi_{\la}(\g,\h)=\phi_{\infty}(\g,\h).
\end{align*}
\end{lemma}}
